All 1,427 maintained tests pass in 289.263 seconds. The 28 focused tests
pass in 13.271 seconds. They check planar admission from raw nullity four,
generic reduction from five to four, complete parent records in the
original coordinate width, legacy replay without the support annotation,
shared-cap boundaries and interruption during source recompilation.
Eleven mutations remove or forge nested Q rays, change the source or allowed
support, change the nested phase or status, add recursive nesting, or alter
parent coverage and Euler data. They are all rejected with the producer,
Q iterator and planar producer disabled during independent replay.

The source audit repeats all 188 retained higher-nullity sectors against
release \code{c21e78bcd}. It completes in 53.634 seconds. Every complete
status agrees. For every source sector, complete standard rays from the
recompiled support equal the retained original-support source oracle in
actual full coordinates. Sixty-five sectors have forced-zero types: 22
fall from matching nullity four to three, 41 stay at four, and two fall from
five to four. Fresh Regina standard enumerations for all three source
triangulations involved reproduce the complete compatible ray lists in all
65 cases. These are actual finite source triangulations, not manufactured
post-elimination kernels.

Automatic discovery recompiles 61 of those sectors and returns 61 augmented
negative standard proofs. The other four return a verified essential Q
witness before recompilation is needed. All 22 raw-nullity-four sectors with
feasible nullity three enter the planar standard producer. Every augmented
parent proof retains the original allowed type list and coordinate width.
The unchanged original full-sector verifier also accepts representative
unannotated proofs at effective dimensions three and four. Independent
publication review replays all 61 augmented proofs with the producers
disabled, and also replays all 31 saved native diagram witnesses from the
preceding release.

The paired benchmark freezes both the preceding discovery module and its
independent checker. Each call constructs its source, completes standard
discovery and all configured work, replays its proof using that arm's
checker, and serializes the full result. Five randomized paired rounds
with identical-baseline A/A controls complete 100 measured calls and twenty
warmups. Both arms agree on every status; every proof is replayed and each
arm's certificate digest is stable across its repetitions.

\begin{center}
\small
\begin{tabular}{lrrrrr}
\toprule
Source & previous ms & current ms & old/new & old A/A & new A/A\\
\midrule
\code{cap\_3\_5\_2-4-to-3} & 164.057 & 41.699 & 3.863 & 1.045 & 1.005\\
\code{cap\_3\_5\_3-4-to-4} & 78.158 & 78.027 & 0.984 & 1.019 & 1.005\\
\code{cap\_3\_5\_3-5-to-4} & 3452.408 & 143.699 & 23.862 & 1.001 & 0.998\\
\code{cap\_1\_2\_3-4-to-4} & 40.810 & 40.753 & 1.004 & 1.002 & 1.002\\
\code{fibonacci\_lst\_09-4-to-4} & 60.635 & 61.058 & 0.999 & 0.988 & 0.988\\
\bottomrule
\end{tabular}
\end{center}


The planar-admission case has paired old/new ratio 3.863, with 684 versus
nine begun candidates, including its completed Q phase. Its incumbent A/A
ratio is 1.045 and the treatment A/A ratio is 1.005; the gain is substantially
larger than those fluctuations. The nullity-five-to-four case has ratio
23.862, with 12,655 versus 369 candidates and A/A controls about 1.001 and
0.998. Both the producer and independent replay avoid their respective
original larger standard arrangements.

The forced-zero case that stays at dimension four is near parity at 0.984;
its candidate count remains 124 in both arms. The full-support positive
control and nonpositive Q control have ratios 1.004 and 0.999. Removing
coordinates without changing effective dimension does not supply a uniform
speedup, and the larger nested proof and source recompilation are included
in the treatment cost. No source-query gain is attributed to improved
coverage of arbitrary knot diagrams.

The runtime release is \code{936e83365}. Tests, the complete source audit,
fresh Regina controls and paired timings are retained under
\code{data/feasible-span-*}. From \code{fast/}, use the module
\code{normal\_orbit\_research.feasible\_span} with its \code{audit} and
\code{benchmark} modes; the audit accepts \code{--fresh-regina}.
Publication checks verify all 609 runtime pins, 607 preceding-release
pins, every timing outcome, the complete source-sector verdicts and all
producer-disabled proof replays. The ray-geometry improvement and feasible
span theorem are local results. A general quasi-polynomial recognizer and
an implemented polynomial support producer remain separate unresolved
requirements. The retained LP formulation gives a route to the latter; this
Q-list integration does not implement that solver.
